% fig_architecture.tex  --  CNN-DDQN+MD Q-network architecture (SCI quality)
% Usage: % fig_architecture.tex  --  CNN-DDQN+MD Q-network architecture (SCI quality)
% Usage: % fig_architecture.tex  --  CNN-DDQN+MD Q-network architecture (SCI quality)
% Usage: % fig_architecture.tex  --  CNN-DDQN+MD Q-network architecture (SCI quality)
% Usage: \input{figures/fig_architecture.tex}
% Requires: \usepackage{tikz}
%           \usetikzlibrary{positioning, arrows.meta, fit, calc, backgrounds}

\begin{figure*}[!htbp]
\centering
%
% ---- Semantic colour palette ----
\definecolor{Cconv}{HTML}{4A90D9}   % CNN conv layers
\definecolor{Cmha}{HTML}{26A69A}    % spatial attention
\definecolor{Cfc}{HTML}{66BB6A}     % fully connected
\definecolor{Ccat}{HTML}{FFD54F}    % concatenation
\definecolor{Cval}{HTML}{EF5350}    % value stream
\definecolor{Cadv}{HTML}{43A047}    % advantage stream
\definecolor{Cqagg}{HTML}{FF9800}   % Q aggregation
\definecolor{Cmask}{HTML}{C62828}   % action mask
\definecolor{Cout}{HTML}{7E57C2}    % output
\definecolor{Cocc}{HTML}{90A4AE}    % occupancy channel
\definecolor{Cgdf}{HTML}{FFB74D}    % goal-distance channel
\definecolor{Cedt}{HTML}{42A5F5}    % EDT channel
%
\resizebox{\textwidth}{!}{%
\begin{tikzpicture}[
    >=Stealth,
    arr/.style  = {->, thick, black!60},
    dlb/.style  = {font=\scriptsize, text=black!55, align=center},
    glb/.style  = {font=\footnotesize\itshape},
]
%
% ==================== 3-D CUBOID MACRO ====================
%   \cub{name}{x}{y}{height}{front-width}{depth}{colour}
%   height → spatial dimension  (shrinks through layers)
%   depth  → channel count      (grows through layers)
\def\cub#1#2#3#4#5#6#7{%
    \pgfmathsetmacro{\cubH}{(#4)/2}%
    \pgfmathsetmacro{\cubW}{(#5)/2}%
    \pgfmathsetmacro{\cubDx}{(#6)*0.28}%
    \pgfmathsetmacro{\cubDy}{(#6)*0.16}%
    % --- right-side face (darkened) ---
    \fill[#7!55!black, draw=black!30, line width=0.3pt]
        ({#2+\cubW},{#3-\cubH})
        -- ({#2+\cubW+\cubDx},{#3-\cubH+\cubDy})
        -- ({#2+\cubW+\cubDx},{#3+\cubH+\cubDy})
        -- ({#2+\cubW},{#3+\cubH}) -- cycle;
    % --- top face (lightened) ---
    \fill[#7!15!white, draw=black!30, line width=0.3pt]
        ({#2-\cubW},{#3+\cubH})
        -- ({#2-\cubW+\cubDx},{#3+\cubH+\cubDy})
        -- ({#2+\cubW+\cubDx},{#3+\cubH+\cubDy})
        -- ({#2+\cubW},{#3+\cubH}) -- cycle;
    % --- front face ---
    \fill[#7!85, draw=black!30, line width=0.3pt]
        ({#2-\cubW},{#3-\cubH}) rectangle ({#2+\cubW},{#3+\cubH});
    % --- named anchors ---
    \coordinate (#1)    at (#2, #3);
    \coordinate (#1-l)  at ({#2-\cubW},          {#3});
    \coordinate (#1-r)  at ({#2+\cubW+\cubDx},   {#3+\cubDy/2});
    \coordinate (#1-t)  at ({#2+\cubDx/2},        {#3+\cubH+\cubDy/2});
    \coordinate (#1-b)  at ({#2},                 {#3-\cubH});
    \coordinate (#1-bl) at ({#2-\cubW},           {#3-\cubH});
    \coordinate (#1-tr) at ({#2+\cubW+\cubDx},   {#3+\cubH+\cubDy});
}
%
% ==================== LAYOUT CONSTANTS ====================
\def\cnnY{1.8}     % CNN branch centre
\def\scaY{-0.6}    % scalar branch
\def\midY{0.6}     % concat / shared-FC / Q-agg
\def\valY{2.3}     % value stream
\def\advY{-1.1}    % advantage stream
%
% ==================== INPUT: 2 MAP CHANNELS ===============
%   back → front: GDF (amber), Occ (grey)
\cub{gdf}{0.42}{\cnnY}{2.4}{0.06}{0.10}{Cgdf}
\cub{occ}{0.22}{\cnnY}{2.4}{0.06}{0.10}{Cocc}
%
\node[dlb, above=0pt] at (occ-t) {O};
\node[dlb, above=0pt] at (gdf-t) {G};
\node[dlb, below=5pt] at (0.32, {\cnnY-1.2})
    {$2{\times}12{\times}12$};
%
% ==================== INPUT: SCALARS ======================
\node[draw, rounded corners=2pt, fill=gray!8, thick,
      minimum width=1.2cm, minimum height=0.50cm,
      align=center, font=\footnotesize]
    (sca) at (0.32, \scaY)
    {Scalars\\[-2pt]{\scriptsize 11-d}};
%
%
% ==================== CNN BACKBONE ========================
%
% Conv1 — 2→32, 3×3, s=1 → 32×12×12
\cub{c1}{2.6}{\cnnY}{2.4}{0.20}{0.55}{Cconv}
\node[font=\footnotesize\bfseries, text=white] at (c1) {32};
\node[dlb, below=5pt] at (c1-b) {Conv$_1$\\$3{\times}3$};
%
% Conv2 — 32→64, 3×3, s=2 → 64×6×6
\cub{c2}{4.6}{\cnnY}{1.5}{0.20}{0.85}{Cconv}
\node[font=\footnotesize\bfseries, text=white] at (c2) {64};
\node[dlb, below=5pt] at (c2-b) {Conv$_2$\\$3{\times}3,\;s{=}2$};
%
% Conv3 — 64→64, 3×3, s=2 → 64×3×3
\cub{c3}{6.4}{\cnnY}{0.9}{0.20}{0.85}{Cconv}
\node[font=\scriptsize\bfseries, text=white] at (c3) {64};
\node[dlb, below=5pt] at (c3-b) {Conv$_3$\\$3{\times}3,\;s{=}2$};
%
% Spatial MHA — 4 heads
\cub{mha}{8.0}{\cnnY}{0.9}{0.20}{0.85}{Cmha}
\node[font=\scriptsize\bfseries, text=white] at (mha) {4h};
\node[dlb, below=5pt] at (mha-b) {Spatial\\[-1pt]MHA};
%
% ---- dimension annotations along the top ----
\node[font=\tiny, text=black!40, above=1pt] at (c1-t) {$12^2$};
\node[font=\tiny, text=black!40, above=1pt] at (c2-t) {$6^2$};
\node[font=\tiny, text=black!40, above=1pt] at (c3-t) {$3^2$};
%
% ---- CNN arrows ----
\draw[arr] (gdf-r) -- (c1-l);
\draw[arr] (c1-r)  -- (c2-l);
\draw[arr] (c2-r)  -- (c3-l);
\draw[arr] (c3-r)  -- (mha-l);
%
%
% ==================== FLATTEN =============================
\node[draw, rounded corners=2pt, fill=Cconv!10, thick,
      minimum height=0.50cm, minimum width=0.65cm,
      font=\footnotesize, align=center]
    (flat) at (9.6, \cnnY) {Flat};
\node[dlb, below=2pt] at (flat.south) {576};
\draw[arr] (mha-r) -- (flat.west);
%
%
% ==================== CONCATENATION =======================
\node[draw, circle, fill=Ccat!35, thick,
      minimum size=0.55cm, inner sep=0pt, font=\normalsize]
    (cat) at (11.2, \midY) {$\oplus$};
\node[dlb, below=2pt] at (cat.south) {587};
%
%   flat (above) ↘ cat ;  sca (below) ↗ cat
\draw[arr, rounded corners=4pt] (flat.east) -| (cat);
\draw[arr, rounded corners=4pt] (sca.east)  -| (cat);
%
%
% ==================== SHARED FC TRUNK =====================
\node[draw, rounded corners=2pt, fill=Cfc!18, thick,
      minimum height=0.55cm, minimum width=1.3cm,
      font=\footnotesize, align=center]
    (sfc) at (13.0, \midY)
    {FC(256)+ReLU};
\draw[arr] (cat) -- (sfc);
%
%
% ==================== DUELING HEAD ========================
%
% ---- fork point ----
\coordinate (fork) at ($(sfc.east)+(0.4,0)$);
%
% ---- Value stream ----
\node[draw, rounded corners=3pt, fill=Cval!15, thick,
      minimum width=1.6cm, minimum height=0.65cm,
      font=\footnotesize, align=center]
    (val) at (16.0, \valY)
    {Value\\[-2pt]{\scriptsize $V(s)\!\to\!\mathbb{R}^1$}};
%
% ---- Advantage stream ----
\node[draw, rounded corners=3pt, fill=Cadv!15, thick,
      minimum width=1.6cm, minimum height=0.65cm,
      font=\footnotesize, align=center]
    (adv) at (16.0, \advY)
    {Advantage\\[-2pt]{\scriptsize $A(s,a)\!\to\!\mathbb{R}^{35}$}};
%
% ---- Y-fork arrows ----
\draw[arr, rounded corners=4pt] (sfc.east) -- (fork) |- (val.west);
\draw[arr, rounded corners=4pt] (fork) |- (adv.west);
%
%
% ==================== Q AGGREGATION =======================
\node[draw, rounded corners=3pt, fill=Cqagg!20, thick,
      minimum width=1.6cm, minimum height=0.65cm,
      font=\footnotesize, align=center]
    (qagg) at (18.8, \midY)
    {$Q{=}V{+}A{-}\overline{A}$\\[-2pt]{\scriptsize 35\;Q-values}};
%
\draw[arr, rounded corners=4pt]
    (val.east) -- ++(0.3,0) |- ($(qagg.west)+(0, 0.08)$);
\draw[arr, rounded corners=4pt]
    (adv.east) -- ++(0.3,0) |- ($(qagg.west)+(0,-0.08)$);
%
%
% ==================== ACTION MASK (7×5 GRID) ==============
\def\gX{20.8}      % grid left x
\def\gY{-0.10}     % grid bottom y  (centred ≈ \midY)
\def\gS{0.22}      % cell size
%
% ---- all cells (valid = green) ----
\foreach \r in {0,...,6}{
    \foreach \c in {0,...,4}{
        \fill[Cadv!15, draw=black!15, line width=0.15pt]
            ({\gX+\c*\gS}, {\gY+\r*\gS})
            rectangle ++({\gS-0.02}, {\gS-0.02});
    }
}
%
% ---- masked cells (red background + × mark) ----
\foreach \r/\c in {0/0,0/1,0/3,0/4, 1/0,1/4, 5/0,5/4, 6/0,6/1,6/3,6/4}{
    \fill[Cmask!12, draw=black!15, line width=0.15pt]
        ({\gX+\c*\gS}, {\gY+\r*\gS})
        rectangle ++({\gS-0.02}, {\gS-0.02});
    \draw[Cmask, line width=0.35pt]
        ({\gX+\c*\gS+0.04}, {\gY+\r*\gS+0.04})
        -- ++({\gS-0.10}, {\gS-0.10});
    \draw[Cmask, line width=0.35pt]
        ({\gX+\c*\gS+\gS-0.06}, {\gY+\r*\gS+0.04})
        -- ++({-\gS+0.10}, {\gS-0.10});
}
%
% ---- grid axis labels ----
\node[dlb, below=4pt] at ({\gX+2.5*\gS}, \gY) {$5\;a$};
\node[dlb, rotate=90]  at ({\gX-0.16}, {\gY+3.5*\gS})
    {$7\;\dot{\delta}$};
%
% ---- kinematic mask input ----
\node[draw, rounded corners=2pt, fill=Cmask!8, thick,
      minimum width=0.9cm, minimum height=0.40cm,
      font=\scriptsize, align=center]
    (kin) at ({\gX+2.5*\gS}, {\gY-0.80})
    {Kin.\ Mask\\[-2pt]{\tiny bool\;35-d}};
\draw[arr] (kin) -- ({\gX+2.5*\gS}, {\gY-0.02});
%
% ---- Q → grid ----
\draw[arr] (qagg.east) -- ({\gX-0.05}, {\gY+3.5*\gS});
%
%
% ==================== OUTPUT ==============================
\node[draw, rounded corners=3pt, fill=Cout!12, thick,
      minimum width=0.7cm, minimum height=0.55cm,
      font=\footnotesize, align=center]
    (out) at (23.5, \midY)
    {$a^*$\\[-2pt]{\scriptsize $(\dot{\delta},\,a)$}};
%
\draw[arr] ({\gX+5*\gS}, {\gY+3.5*\gS})
    -- node[above, font=\scriptsize, text=black!50] {argmax}
    (out.west);
%
%
% ==================== BACKGROUND GROUPS ===================
\begin{scope}[on background layer]
    %
    % CNN Feature Extractor
    \node[draw=Cconv!40, dashed, rounded corners=5pt,
          inner sep=7pt, fill=Cconv!3,
          fit=(c1-bl)(c1-tr)(mha-tr)(flat),
          label={[glb, Cconv!60]above:CNN Feature Extractor}] {};
    %
    % Dueling Head (MD)
    \node[draw=Cqagg!50, dashed, rounded corners=5pt,
          inner sep=7pt, fill=Cqagg!3,
          fit=(val)(adv)(qagg),
          label={[glb, Cqagg!70]above:Dueling Head (MD)}] {};
    %
    % Action Masking
    \coordinate (mask-bl) at ({\gX-0.12}, {\gY-0.12});
    \coordinate (mask-tr) at ({\gX+5*\gS+0.10}, {\gY+7*\gS+0.10});
    \node[draw=Cmask!35, dashed, rounded corners=5pt,
          inner sep=7pt, fill=Cmask!2,
          fit=(mask-bl)(mask-tr)(kin)(out),
          label={[glb, Cmask!50]above:Action Masking}] {};
    %
\end{scope}
%
\end{tikzpicture}%
}% end resizebox
%
\caption{%
Architecture of the CNN-DDQN+MD Q-network.
A convolutional branch extracts spatial features from two
$12{\times}12$ map channels
(occupancy, goal distance field)
through three convolution layers with ReLU activations
and spatial multi-head attention (MHA),
yielding a 576-dimensional feature vector.
This vector is concatenated with 11 scalar state features
(pose, velocity, steering angle, goal bearing)
and processed by a shared FC layer before splitting
into a value stream~$V(s)$ and an advantage stream~$A(s,a)$
following the dueling architecture.
The head produces 35 joint Q-values
($7$ steering rates\,$\times$\,$5$ accelerations),
and a kinematic admissibility mask zeroes infeasible actions
before argmax policy extraction.}
\label{fig:architecture}
\end{figure*}

% Requires: \usepackage{tikz}
%           \usetikzlibrary{positioning, arrows.meta, fit, calc, backgrounds}

\begin{figure*}[!htbp]
\centering
%
% ---- Semantic colour palette ----
\definecolor{Cconv}{HTML}{4A90D9}   % CNN conv layers
\definecolor{Cmha}{HTML}{26A69A}    % spatial attention
\definecolor{Cfc}{HTML}{66BB6A}     % fully connected
\definecolor{Ccat}{HTML}{FFD54F}    % concatenation
\definecolor{Cval}{HTML}{EF5350}    % value stream
\definecolor{Cadv}{HTML}{43A047}    % advantage stream
\definecolor{Cqagg}{HTML}{FF9800}   % Q aggregation
\definecolor{Cmask}{HTML}{C62828}   % action mask
\definecolor{Cout}{HTML}{7E57C2}    % output
\definecolor{Cocc}{HTML}{90A4AE}    % occupancy channel
\definecolor{Cgdf}{HTML}{FFB74D}    % goal-distance channel
\definecolor{Cedt}{HTML}{42A5F5}    % EDT channel
%
\resizebox{\textwidth}{!}{%
\begin{tikzpicture}[
    >=Stealth,
    arr/.style  = {->, thick, black!60},
    dlb/.style  = {font=\scriptsize, text=black!55, align=center},
    glb/.style  = {font=\footnotesize\itshape},
]
%
% ==================== 3-D CUBOID MACRO ====================
%   \cub{name}{x}{y}{height}{front-width}{depth}{colour}
%   height → spatial dimension  (shrinks through layers)
%   depth  → channel count      (grows through layers)
\def\cub#1#2#3#4#5#6#7{%
    \pgfmathsetmacro{\cubH}{(#4)/2}%
    \pgfmathsetmacro{\cubW}{(#5)/2}%
    \pgfmathsetmacro{\cubDx}{(#6)*0.28}%
    \pgfmathsetmacro{\cubDy}{(#6)*0.16}%
    % --- right-side face (darkened) ---
    \fill[#7!55!black, draw=black!30, line width=0.3pt]
        ({#2+\cubW},{#3-\cubH})
        -- ({#2+\cubW+\cubDx},{#3-\cubH+\cubDy})
        -- ({#2+\cubW+\cubDx},{#3+\cubH+\cubDy})
        -- ({#2+\cubW},{#3+\cubH}) -- cycle;
    % --- top face (lightened) ---
    \fill[#7!15!white, draw=black!30, line width=0.3pt]
        ({#2-\cubW},{#3+\cubH})
        -- ({#2-\cubW+\cubDx},{#3+\cubH+\cubDy})
        -- ({#2+\cubW+\cubDx},{#3+\cubH+\cubDy})
        -- ({#2+\cubW},{#3+\cubH}) -- cycle;
    % --- front face ---
    \fill[#7!85, draw=black!30, line width=0.3pt]
        ({#2-\cubW},{#3-\cubH}) rectangle ({#2+\cubW},{#3+\cubH});
    % --- named anchors ---
    \coordinate (#1)    at (#2, #3);
    \coordinate (#1-l)  at ({#2-\cubW},          {#3});
    \coordinate (#1-r)  at ({#2+\cubW+\cubDx},   {#3+\cubDy/2});
    \coordinate (#1-t)  at ({#2+\cubDx/2},        {#3+\cubH+\cubDy/2});
    \coordinate (#1-b)  at ({#2},                 {#3-\cubH});
    \coordinate (#1-bl) at ({#2-\cubW},           {#3-\cubH});
    \coordinate (#1-tr) at ({#2+\cubW+\cubDx},   {#3+\cubH+\cubDy});
}
%
% ==================== LAYOUT CONSTANTS ====================
\def\cnnY{1.8}     % CNN branch centre
\def\scaY{-0.6}    % scalar branch
\def\midY{0.6}     % concat / shared-FC / Q-agg
\def\valY{2.3}     % value stream
\def\advY{-1.1}    % advantage stream
%
% ==================== INPUT: 2 MAP CHANNELS ===============
%   back → front: GDF (amber), Occ (grey)
\cub{gdf}{0.42}{\cnnY}{2.4}{0.06}{0.10}{Cgdf}
\cub{occ}{0.22}{\cnnY}{2.4}{0.06}{0.10}{Cocc}
%
\node[dlb, above=0pt] at (occ-t) {O};
\node[dlb, above=0pt] at (gdf-t) {G};
\node[dlb, below=5pt] at (0.32, {\cnnY-1.2})
    {$2{\times}12{\times}12$};
%
% ==================== INPUT: SCALARS ======================
\node[draw, rounded corners=2pt, fill=gray!8, thick,
      minimum width=1.2cm, minimum height=0.50cm,
      align=center, font=\footnotesize]
    (sca) at (0.32, \scaY)
    {Scalars\\[-2pt]{\scriptsize 11-d}};
%
%
% ==================== CNN BACKBONE ========================
%
% Conv1 — 2→32, 3×3, s=1 → 32×12×12
\cub{c1}{2.6}{\cnnY}{2.4}{0.20}{0.55}{Cconv}
\node[font=\footnotesize\bfseries, text=white] at (c1) {32};
\node[dlb, below=5pt] at (c1-b) {Conv$_1$\\$3{\times}3$};
%
% Conv2 — 32→64, 3×3, s=2 → 64×6×6
\cub{c2}{4.6}{\cnnY}{1.5}{0.20}{0.85}{Cconv}
\node[font=\footnotesize\bfseries, text=white] at (c2) {64};
\node[dlb, below=5pt] at (c2-b) {Conv$_2$\\$3{\times}3,\;s{=}2$};
%
% Conv3 — 64→64, 3×3, s=2 → 64×3×3
\cub{c3}{6.4}{\cnnY}{0.9}{0.20}{0.85}{Cconv}
\node[font=\scriptsize\bfseries, text=white] at (c3) {64};
\node[dlb, below=5pt] at (c3-b) {Conv$_3$\\$3{\times}3,\;s{=}2$};
%
% Spatial MHA — 4 heads
\cub{mha}{8.0}{\cnnY}{0.9}{0.20}{0.85}{Cmha}
\node[font=\scriptsize\bfseries, text=white] at (mha) {4h};
\node[dlb, below=5pt] at (mha-b) {Spatial\\[-1pt]MHA};
%
% ---- dimension annotations along the top ----
\node[font=\tiny, text=black!40, above=1pt] at (c1-t) {$12^2$};
\node[font=\tiny, text=black!40, above=1pt] at (c2-t) {$6^2$};
\node[font=\tiny, text=black!40, above=1pt] at (c3-t) {$3^2$};
%
% ---- CNN arrows ----
\draw[arr] (gdf-r) -- (c1-l);
\draw[arr] (c1-r)  -- (c2-l);
\draw[arr] (c2-r)  -- (c3-l);
\draw[arr] (c3-r)  -- (mha-l);
%
%
% ==================== FLATTEN =============================
\node[draw, rounded corners=2pt, fill=Cconv!10, thick,
      minimum height=0.50cm, minimum width=0.65cm,
      font=\footnotesize, align=center]
    (flat) at (9.6, \cnnY) {Flat};
\node[dlb, below=2pt] at (flat.south) {576};
\draw[arr] (mha-r) -- (flat.west);
%
%
% ==================== CONCATENATION =======================
\node[draw, circle, fill=Ccat!35, thick,
      minimum size=0.55cm, inner sep=0pt, font=\normalsize]
    (cat) at (11.2, \midY) {$\oplus$};
\node[dlb, below=2pt] at (cat.south) {587};
%
%   flat (above) ↘ cat ;  sca (below) ↗ cat
\draw[arr, rounded corners=4pt] (flat.east) -| (cat);
\draw[arr, rounded corners=4pt] (sca.east)  -| (cat);
%
%
% ==================== SHARED FC TRUNK =====================
\node[draw, rounded corners=2pt, fill=Cfc!18, thick,
      minimum height=0.55cm, minimum width=1.3cm,
      font=\footnotesize, align=center]
    (sfc) at (13.0, \midY)
    {FC(256)+ReLU};
\draw[arr] (cat) -- (sfc);
%
%
% ==================== DUELING HEAD ========================
%
% ---- fork point ----
\coordinate (fork) at ($(sfc.east)+(0.4,0)$);
%
% ---- Value stream ----
\node[draw, rounded corners=3pt, fill=Cval!15, thick,
      minimum width=1.6cm, minimum height=0.65cm,
      font=\footnotesize, align=center]
    (val) at (16.0, \valY)
    {Value\\[-2pt]{\scriptsize $V(s)\!\to\!\mathbb{R}^1$}};
%
% ---- Advantage stream ----
\node[draw, rounded corners=3pt, fill=Cadv!15, thick,
      minimum width=1.6cm, minimum height=0.65cm,
      font=\footnotesize, align=center]
    (adv) at (16.0, \advY)
    {Advantage\\[-2pt]{\scriptsize $A(s,a)\!\to\!\mathbb{R}^{35}$}};
%
% ---- Y-fork arrows ----
\draw[arr, rounded corners=4pt] (sfc.east) -- (fork) |- (val.west);
\draw[arr, rounded corners=4pt] (fork) |- (adv.west);
%
%
% ==================== Q AGGREGATION =======================
\node[draw, rounded corners=3pt, fill=Cqagg!20, thick,
      minimum width=1.6cm, minimum height=0.65cm,
      font=\footnotesize, align=center]
    (qagg) at (18.8, \midY)
    {$Q{=}V{+}A{-}\overline{A}$\\[-2pt]{\scriptsize 35\;Q-values}};
%
\draw[arr, rounded corners=4pt]
    (val.east) -- ++(0.3,0) |- ($(qagg.west)+(0, 0.08)$);
\draw[arr, rounded corners=4pt]
    (adv.east) -- ++(0.3,0) |- ($(qagg.west)+(0,-0.08)$);
%
%
% ==================== ACTION MASK (7×5 GRID) ==============
\def\gX{20.8}      % grid left x
\def\gY{-0.10}     % grid bottom y  (centred ≈ \midY)
\def\gS{0.22}      % cell size
%
% ---- all cells (valid = green) ----
\foreach \r in {0,...,6}{
    \foreach \c in {0,...,4}{
        \fill[Cadv!15, draw=black!15, line width=0.15pt]
            ({\gX+\c*\gS}, {\gY+\r*\gS})
            rectangle ++({\gS-0.02}, {\gS-0.02});
    }
}
%
% ---- masked cells (red background + × mark) ----
\foreach \r/\c in {0/0,0/1,0/3,0/4, 1/0,1/4, 5/0,5/4, 6/0,6/1,6/3,6/4}{
    \fill[Cmask!12, draw=black!15, line width=0.15pt]
        ({\gX+\c*\gS}, {\gY+\r*\gS})
        rectangle ++({\gS-0.02}, {\gS-0.02});
    \draw[Cmask, line width=0.35pt]
        ({\gX+\c*\gS+0.04}, {\gY+\r*\gS+0.04})
        -- ++({\gS-0.10}, {\gS-0.10});
    \draw[Cmask, line width=0.35pt]
        ({\gX+\c*\gS+\gS-0.06}, {\gY+\r*\gS+0.04})
        -- ++({-\gS+0.10}, {\gS-0.10});
}
%
% ---- grid axis labels ----
\node[dlb, below=4pt] at ({\gX+2.5*\gS}, \gY) {$5\;a$};
\node[dlb, rotate=90]  at ({\gX-0.16}, {\gY+3.5*\gS})
    {$7\;\dot{\delta}$};
%
% ---- kinematic mask input ----
\node[draw, rounded corners=2pt, fill=Cmask!8, thick,
      minimum width=0.9cm, minimum height=0.40cm,
      font=\scriptsize, align=center]
    (kin) at ({\gX+2.5*\gS}, {\gY-0.80})
    {Kin.\ Mask\\[-2pt]{\tiny bool\;35-d}};
\draw[arr] (kin) -- ({\gX+2.5*\gS}, {\gY-0.02});
%
% ---- Q → grid ----
\draw[arr] (qagg.east) -- ({\gX-0.05}, {\gY+3.5*\gS});
%
%
% ==================== OUTPUT ==============================
\node[draw, rounded corners=3pt, fill=Cout!12, thick,
      minimum width=0.7cm, minimum height=0.55cm,
      font=\footnotesize, align=center]
    (out) at (23.5, \midY)
    {$a^*$\\[-2pt]{\scriptsize $(\dot{\delta},\,a)$}};
%
\draw[arr] ({\gX+5*\gS}, {\gY+3.5*\gS})
    -- node[above, font=\scriptsize, text=black!50] {argmax}
    (out.west);
%
%
% ==================== BACKGROUND GROUPS ===================
\begin{scope}[on background layer]
    %
    % CNN Feature Extractor
    \node[draw=Cconv!40, dashed, rounded corners=5pt,
          inner sep=7pt, fill=Cconv!3,
          fit=(c1-bl)(c1-tr)(mha-tr)(flat),
          label={[glb, Cconv!60]above:CNN Feature Extractor}] {};
    %
    % Dueling Head (MD)
    \node[draw=Cqagg!50, dashed, rounded corners=5pt,
          inner sep=7pt, fill=Cqagg!3,
          fit=(val)(adv)(qagg),
          label={[glb, Cqagg!70]above:Dueling Head (MD)}] {};
    %
    % Action Masking
    \coordinate (mask-bl) at ({\gX-0.12}, {\gY-0.12});
    \coordinate (mask-tr) at ({\gX+5*\gS+0.10}, {\gY+7*\gS+0.10});
    \node[draw=Cmask!35, dashed, rounded corners=5pt,
          inner sep=7pt, fill=Cmask!2,
          fit=(mask-bl)(mask-tr)(kin)(out),
          label={[glb, Cmask!50]above:Action Masking}] {};
    %
\end{scope}
%
\end{tikzpicture}%
}% end resizebox
%
\caption{%
Architecture of the CNN-DDQN+MD Q-network.
A convolutional branch extracts spatial features from two
$12{\times}12$ map channels
(occupancy, goal distance field)
through three convolution layers with ReLU activations
and spatial multi-head attention (MHA),
yielding a 576-dimensional feature vector.
This vector is concatenated with 11 scalar state features
(pose, velocity, steering angle, goal bearing)
and processed by a shared FC layer before splitting
into a value stream~$V(s)$ and an advantage stream~$A(s,a)$
following the dueling architecture.
The head produces 35 joint Q-values
($7$ steering rates\,$\times$\,$5$ accelerations),
and a kinematic admissibility mask zeroes infeasible actions
before argmax policy extraction.}
\label{fig:architecture}
\end{figure*}

% Requires: \usepackage{tikz}
%           \usetikzlibrary{positioning, arrows.meta, fit, calc, backgrounds}

\begin{figure*}[!htbp]
\centering
%
% ---- Semantic colour palette ----
\definecolor{Cconv}{HTML}{4A90D9}   % CNN conv layers
\definecolor{Cmha}{HTML}{26A69A}    % spatial attention
\definecolor{Cfc}{HTML}{66BB6A}     % fully connected
\definecolor{Ccat}{HTML}{FFD54F}    % concatenation
\definecolor{Cval}{HTML}{EF5350}    % value stream
\definecolor{Cadv}{HTML}{43A047}    % advantage stream
\definecolor{Cqagg}{HTML}{FF9800}   % Q aggregation
\definecolor{Cmask}{HTML}{C62828}   % action mask
\definecolor{Cout}{HTML}{7E57C2}    % output
\definecolor{Cocc}{HTML}{90A4AE}    % occupancy channel
\definecolor{Cgdf}{HTML}{FFB74D}    % goal-distance channel
\definecolor{Cedt}{HTML}{42A5F5}    % EDT channel
%
\resizebox{\textwidth}{!}{%
\begin{tikzpicture}[
    >=Stealth,
    arr/.style  = {->, thick, black!60},
    dlb/.style  = {font=\scriptsize, text=black!55, align=center},
    glb/.style  = {font=\footnotesize\itshape},
]
%
% ==================== 3-D CUBOID MACRO ====================
%   \cub{name}{x}{y}{height}{front-width}{depth}{colour}
%   height → spatial dimension  (shrinks through layers)
%   depth  → channel count      (grows through layers)
\def\cub#1#2#3#4#5#6#7{%
    \pgfmathsetmacro{\cubH}{(#4)/2}%
    \pgfmathsetmacro{\cubW}{(#5)/2}%
    \pgfmathsetmacro{\cubDx}{(#6)*0.28}%
    \pgfmathsetmacro{\cubDy}{(#6)*0.16}%
    % --- right-side face (darkened) ---
    \fill[#7!55!black, draw=black!30, line width=0.3pt]
        ({#2+\cubW},{#3-\cubH})
        -- ({#2+\cubW+\cubDx},{#3-\cubH+\cubDy})
        -- ({#2+\cubW+\cubDx},{#3+\cubH+\cubDy})
        -- ({#2+\cubW},{#3+\cubH}) -- cycle;
    % --- top face (lightened) ---
    \fill[#7!15!white, draw=black!30, line width=0.3pt]
        ({#2-\cubW},{#3+\cubH})
        -- ({#2-\cubW+\cubDx},{#3+\cubH+\cubDy})
        -- ({#2+\cubW+\cubDx},{#3+\cubH+\cubDy})
        -- ({#2+\cubW},{#3+\cubH}) -- cycle;
    % --- front face ---
    \fill[#7!85, draw=black!30, line width=0.3pt]
        ({#2-\cubW},{#3-\cubH}) rectangle ({#2+\cubW},{#3+\cubH});
    % --- named anchors ---
    \coordinate (#1)    at (#2, #3);
    \coordinate (#1-l)  at ({#2-\cubW},          {#3});
    \coordinate (#1-r)  at ({#2+\cubW+\cubDx},   {#3+\cubDy/2});
    \coordinate (#1-t)  at ({#2+\cubDx/2},        {#3+\cubH+\cubDy/2});
    \coordinate (#1-b)  at ({#2},                 {#3-\cubH});
    \coordinate (#1-bl) at ({#2-\cubW},           {#3-\cubH});
    \coordinate (#1-tr) at ({#2+\cubW+\cubDx},   {#3+\cubH+\cubDy});
}
%
% ==================== LAYOUT CONSTANTS ====================
\def\cnnY{1.8}     % CNN branch centre
\def\scaY{-0.6}    % scalar branch
\def\midY{0.6}     % concat / shared-FC / Q-agg
\def\valY{2.3}     % value stream
\def\advY{-1.1}    % advantage stream
%
% ==================== INPUT: 2 MAP CHANNELS ===============
%   back → front: GDF (amber), Occ (grey)
\cub{gdf}{0.42}{\cnnY}{2.4}{0.06}{0.10}{Cgdf}
\cub{occ}{0.22}{\cnnY}{2.4}{0.06}{0.10}{Cocc}
%
\node[dlb, above=0pt] at (occ-t) {O};
\node[dlb, above=0pt] at (gdf-t) {G};
\node[dlb, below=5pt] at (0.32, {\cnnY-1.2})
    {$2{\times}12{\times}12$};
%
% ==================== INPUT: SCALARS ======================
\node[draw, rounded corners=2pt, fill=gray!8, thick,
      minimum width=1.2cm, minimum height=0.50cm,
      align=center, font=\footnotesize]
    (sca) at (0.32, \scaY)
    {Scalars\\[-2pt]{\scriptsize 11-d}};
%
%
% ==================== CNN BACKBONE ========================
%
% Conv1 — 2→32, 3×3, s=1 → 32×12×12
\cub{c1}{2.6}{\cnnY}{2.4}{0.20}{0.55}{Cconv}
\node[font=\footnotesize\bfseries, text=white] at (c1) {32};
\node[dlb, below=5pt] at (c1-b) {Conv$_1$\\$3{\times}3$};
%
% Conv2 — 32→64, 3×3, s=2 → 64×6×6
\cub{c2}{4.6}{\cnnY}{1.5}{0.20}{0.85}{Cconv}
\node[font=\footnotesize\bfseries, text=white] at (c2) {64};
\node[dlb, below=5pt] at (c2-b) {Conv$_2$\\$3{\times}3,\;s{=}2$};
%
% Conv3 — 64→64, 3×3, s=2 → 64×3×3
\cub{c3}{6.4}{\cnnY}{0.9}{0.20}{0.85}{Cconv}
\node[font=\scriptsize\bfseries, text=white] at (c3) {64};
\node[dlb, below=5pt] at (c3-b) {Conv$_3$\\$3{\times}3,\;s{=}2$};
%
% Spatial MHA — 4 heads
\cub{mha}{8.0}{\cnnY}{0.9}{0.20}{0.85}{Cmha}
\node[font=\scriptsize\bfseries, text=white] at (mha) {4h};
\node[dlb, below=5pt] at (mha-b) {Spatial\\[-1pt]MHA};
%
% ---- dimension annotations along the top ----
\node[font=\tiny, text=black!40, above=1pt] at (c1-t) {$12^2$};
\node[font=\tiny, text=black!40, above=1pt] at (c2-t) {$6^2$};
\node[font=\tiny, text=black!40, above=1pt] at (c3-t) {$3^2$};
%
% ---- CNN arrows ----
\draw[arr] (gdf-r) -- (c1-l);
\draw[arr] (c1-r)  -- (c2-l);
\draw[arr] (c2-r)  -- (c3-l);
\draw[arr] (c3-r)  -- (mha-l);
%
%
% ==================== FLATTEN =============================
\node[draw, rounded corners=2pt, fill=Cconv!10, thick,
      minimum height=0.50cm, minimum width=0.65cm,
      font=\footnotesize, align=center]
    (flat) at (9.6, \cnnY) {Flat};
\node[dlb, below=2pt] at (flat.south) {576};
\draw[arr] (mha-r) -- (flat.west);
%
%
% ==================== CONCATENATION =======================
\node[draw, circle, fill=Ccat!35, thick,
      minimum size=0.55cm, inner sep=0pt, font=\normalsize]
    (cat) at (11.2, \midY) {$\oplus$};
\node[dlb, below=2pt] at (cat.south) {587};
%
%   flat (above) ↘ cat ;  sca (below) ↗ cat
\draw[arr, rounded corners=4pt] (flat.east) -| (cat);
\draw[arr, rounded corners=4pt] (sca.east)  -| (cat);
%
%
% ==================== SHARED FC TRUNK =====================
\node[draw, rounded corners=2pt, fill=Cfc!18, thick,
      minimum height=0.55cm, minimum width=1.3cm,
      font=\footnotesize, align=center]
    (sfc) at (13.0, \midY)
    {FC(256)+ReLU};
\draw[arr] (cat) -- (sfc);
%
%
% ==================== DUELING HEAD ========================
%
% ---- fork point ----
\coordinate (fork) at ($(sfc.east)+(0.4,0)$);
%
% ---- Value stream ----
\node[draw, rounded corners=3pt, fill=Cval!15, thick,
      minimum width=1.6cm, minimum height=0.65cm,
      font=\footnotesize, align=center]
    (val) at (16.0, \valY)
    {Value\\[-2pt]{\scriptsize $V(s)\!\to\!\mathbb{R}^1$}};
%
% ---- Advantage stream ----
\node[draw, rounded corners=3pt, fill=Cadv!15, thick,
      minimum width=1.6cm, minimum height=0.65cm,
      font=\footnotesize, align=center]
    (adv) at (16.0, \advY)
    {Advantage\\[-2pt]{\scriptsize $A(s,a)\!\to\!\mathbb{R}^{35}$}};
%
% ---- Y-fork arrows ----
\draw[arr, rounded corners=4pt] (sfc.east) -- (fork) |- (val.west);
\draw[arr, rounded corners=4pt] (fork) |- (adv.west);
%
%
% ==================== Q AGGREGATION =======================
\node[draw, rounded corners=3pt, fill=Cqagg!20, thick,
      minimum width=1.6cm, minimum height=0.65cm,
      font=\footnotesize, align=center]
    (qagg) at (18.8, \midY)
    {$Q{=}V{+}A{-}\overline{A}$\\[-2pt]{\scriptsize 35\;Q-values}};
%
\draw[arr, rounded corners=4pt]
    (val.east) -- ++(0.3,0) |- ($(qagg.west)+(0, 0.08)$);
\draw[arr, rounded corners=4pt]
    (adv.east) -- ++(0.3,0) |- ($(qagg.west)+(0,-0.08)$);
%
%
% ==================== ACTION MASK (7×5 GRID) ==============
\def\gX{20.8}      % grid left x
\def\gY{-0.10}     % grid bottom y  (centred ≈ \midY)
\def\gS{0.22}      % cell size
%
% ---- all cells (valid = green) ----
\foreach \r in {0,...,6}{
    \foreach \c in {0,...,4}{
        \fill[Cadv!15, draw=black!15, line width=0.15pt]
            ({\gX+\c*\gS}, {\gY+\r*\gS})
            rectangle ++({\gS-0.02}, {\gS-0.02});
    }
}
%
% ---- masked cells (red background + × mark) ----
\foreach \r/\c in {0/0,0/1,0/3,0/4, 1/0,1/4, 5/0,5/4, 6/0,6/1,6/3,6/4}{
    \fill[Cmask!12, draw=black!15, line width=0.15pt]
        ({\gX+\c*\gS}, {\gY+\r*\gS})
        rectangle ++({\gS-0.02}, {\gS-0.02});
    \draw[Cmask, line width=0.35pt]
        ({\gX+\c*\gS+0.04}, {\gY+\r*\gS+0.04})
        -- ++({\gS-0.10}, {\gS-0.10});
    \draw[Cmask, line width=0.35pt]
        ({\gX+\c*\gS+\gS-0.06}, {\gY+\r*\gS+0.04})
        -- ++({-\gS+0.10}, {\gS-0.10});
}
%
% ---- grid axis labels ----
\node[dlb, below=4pt] at ({\gX+2.5*\gS}, \gY) {$5\;a$};
\node[dlb, rotate=90]  at ({\gX-0.16}, {\gY+3.5*\gS})
    {$7\;\dot{\delta}$};
%
% ---- kinematic mask input ----
\node[draw, rounded corners=2pt, fill=Cmask!8, thick,
      minimum width=0.9cm, minimum height=0.40cm,
      font=\scriptsize, align=center]
    (kin) at ({\gX+2.5*\gS}, {\gY-0.80})
    {Kin.\ Mask\\[-2pt]{\tiny bool\;35-d}};
\draw[arr] (kin) -- ({\gX+2.5*\gS}, {\gY-0.02});
%
% ---- Q → grid ----
\draw[arr] (qagg.east) -- ({\gX-0.05}, {\gY+3.5*\gS});
%
%
% ==================== OUTPUT ==============================
\node[draw, rounded corners=3pt, fill=Cout!12, thick,
      minimum width=0.7cm, minimum height=0.55cm,
      font=\footnotesize, align=center]
    (out) at (23.5, \midY)
    {$a^*$\\[-2pt]{\scriptsize $(\dot{\delta},\,a)$}};
%
\draw[arr] ({\gX+5*\gS}, {\gY+3.5*\gS})
    -- node[above, font=\scriptsize, text=black!50] {argmax}
    (out.west);
%
%
% ==================== BACKGROUND GROUPS ===================
\begin{scope}[on background layer]
    %
    % CNN Feature Extractor
    \node[draw=Cconv!40, dashed, rounded corners=5pt,
          inner sep=7pt, fill=Cconv!3,
          fit=(c1-bl)(c1-tr)(mha-tr)(flat),
          label={[glb, Cconv!60]above:CNN Feature Extractor}] {};
    %
    % Dueling Head (MD)
    \node[draw=Cqagg!50, dashed, rounded corners=5pt,
          inner sep=7pt, fill=Cqagg!3,
          fit=(val)(adv)(qagg),
          label={[glb, Cqagg!70]above:Dueling Head (MD)}] {};
    %
    % Action Masking
    \coordinate (mask-bl) at ({\gX-0.12}, {\gY-0.12});
    \coordinate (mask-tr) at ({\gX+5*\gS+0.10}, {\gY+7*\gS+0.10});
    \node[draw=Cmask!35, dashed, rounded corners=5pt,
          inner sep=7pt, fill=Cmask!2,
          fit=(mask-bl)(mask-tr)(kin)(out),
          label={[glb, Cmask!50]above:Action Masking}] {};
    %
\end{scope}
%
\end{tikzpicture}%
}% end resizebox
%
\caption{%
Architecture of the CNN-DDQN+MD Q-network.
A convolutional branch extracts spatial features from two
$12{\times}12$ map channels
(occupancy, goal distance field)
through three convolution layers with ReLU activations
and spatial multi-head attention (MHA),
yielding a 576-dimensional feature vector.
This vector is concatenated with 11 scalar state features
(pose, velocity, steering angle, goal bearing)
and processed by a shared FC layer before splitting
into a value stream~$V(s)$ and an advantage stream~$A(s,a)$
following the dueling architecture.
The head produces 35 joint Q-values
($7$ steering rates\,$\times$\,$5$ accelerations),
and a kinematic admissibility mask zeroes infeasible actions
before argmax policy extraction.}
\label{fig:architecture}
\end{figure*}

% Requires: \usepackage{tikz}
%           \usetikzlibrary{positioning, arrows.meta, fit, calc, backgrounds}

\begin{figure*}[!htbp]
\centering
%
% ---- Semantic colour palette ----
\definecolor{Cconv}{HTML}{4A90D9}   % CNN conv layers
\definecolor{Cmha}{HTML}{26A69A}    % spatial attention
\definecolor{Cfc}{HTML}{66BB6A}     % fully connected
\definecolor{Ccat}{HTML}{FFD54F}    % concatenation
\definecolor{Cval}{HTML}{EF5350}    % value stream
\definecolor{Cadv}{HTML}{43A047}    % advantage stream
\definecolor{Cqagg}{HTML}{FF9800}   % Q aggregation
\definecolor{Cmask}{HTML}{C62828}   % action mask
\definecolor{Cout}{HTML}{7E57C2}    % output
\definecolor{Cocc}{HTML}{90A4AE}    % occupancy channel
\definecolor{Cgdf}{HTML}{FFB74D}    % goal-distance channel
\definecolor{Cedt}{HTML}{42A5F5}    % EDT channel
%
\resizebox{\textwidth}{!}{%
\begin{tikzpicture}[
    >=Stealth,
    arr/.style  = {->, thick, black!60},
    dlb/.style  = {font=\scriptsize, text=black!55, align=center},
    glb/.style  = {font=\footnotesize\itshape},
]
%
% ==================== 3-D CUBOID MACRO ====================
%   \cub{name}{x}{y}{height}{front-width}{depth}{colour}
%   height → spatial dimension  (shrinks through layers)
%   depth  → channel count      (grows through layers)
\def\cub#1#2#3#4#5#6#7{%
    \pgfmathsetmacro{\cubH}{(#4)/2}%
    \pgfmathsetmacro{\cubW}{(#5)/2}%
    \pgfmathsetmacro{\cubDx}{(#6)*0.28}%
    \pgfmathsetmacro{\cubDy}{(#6)*0.16}%
    % --- right-side face (darkened) ---
    \fill[#7!55!black, draw=black!30, line width=0.3pt]
        ({#2+\cubW},{#3-\cubH})
        -- ({#2+\cubW+\cubDx},{#3-\cubH+\cubDy})
        -- ({#2+\cubW+\cubDx},{#3+\cubH+\cubDy})
        -- ({#2+\cubW},{#3+\cubH}) -- cycle;
    % --- top face (lightened) ---
    \fill[#7!15!white, draw=black!30, line width=0.3pt]
        ({#2-\cubW},{#3+\cubH})
        -- ({#2-\cubW+\cubDx},{#3+\cubH+\cubDy})
        -- ({#2+\cubW+\cubDx},{#3+\cubH+\cubDy})
        -- ({#2+\cubW},{#3+\cubH}) -- cycle;
    % --- front face ---
    \fill[#7!85, draw=black!30, line width=0.3pt]
        ({#2-\cubW},{#3-\cubH}) rectangle ({#2+\cubW},{#3+\cubH});
    % --- named anchors ---
    \coordinate (#1)    at (#2, #3);
    \coordinate (#1-l)  at ({#2-\cubW},          {#3});
    \coordinate (#1-r)  at ({#2+\cubW+\cubDx},   {#3+\cubDy/2});
    \coordinate (#1-t)  at ({#2+\cubDx/2},        {#3+\cubH+\cubDy/2});
    \coordinate (#1-b)  at ({#2},                 {#3-\cubH});
    \coordinate (#1-bl) at ({#2-\cubW},           {#3-\cubH});
    \coordinate (#1-tr) at ({#2+\cubW+\cubDx},   {#3+\cubH+\cubDy});
}
%
% ==================== LAYOUT CONSTANTS ====================
\def\cnnY{1.8}     % CNN branch centre
\def\scaY{-0.6}    % scalar branch
\def\midY{0.6}     % concat / shared-FC / Q-agg
\def\valY{2.3}     % value stream
\def\advY{-1.1}    % advantage stream
%
% ==================== INPUT: 2 MAP CHANNELS ===============
%   back → front: GDF (amber), Occ (grey)
\cub{gdf}{0.42}{\cnnY}{2.4}{0.06}{0.10}{Cgdf}
\cub{occ}{0.22}{\cnnY}{2.4}{0.06}{0.10}{Cocc}
%
\node[dlb, above=0pt] at (occ-t) {O};
\node[dlb, above=0pt] at (gdf-t) {G};
\node[dlb, below=5pt] at (0.32, {\cnnY-1.2})
    {$2{\times}12{\times}12$};
%
% ==================== INPUT: SCALARS ======================
\node[draw, rounded corners=2pt, fill=gray!8, thick,
      minimum width=1.2cm, minimum height=0.50cm,
      align=center, font=\footnotesize]
    (sca) at (0.32, \scaY)
    {Scalars\\[-2pt]{\scriptsize 11-d}};
%
%
% ==================== CNN BACKBONE ========================
%
% Conv1 — 2→32, 3×3, s=1 → 32×12×12
\cub{c1}{2.6}{\cnnY}{2.4}{0.20}{0.55}{Cconv}
\node[font=\footnotesize\bfseries, text=white] at (c1) {32};
\node[dlb, below=5pt] at (c1-b) {Conv$_1$\\$3{\times}3$};
%
% Conv2 — 32→64, 3×3, s=2 → 64×6×6
\cub{c2}{4.6}{\cnnY}{1.5}{0.20}{0.85}{Cconv}
\node[font=\footnotesize\bfseries, text=white] at (c2) {64};
\node[dlb, below=5pt] at (c2-b) {Conv$_2$\\$3{\times}3,\;s{=}2$};
%
% Conv3 — 64→64, 3×3, s=2 → 64×3×3
\cub{c3}{6.4}{\cnnY}{0.9}{0.20}{0.85}{Cconv}
\node[font=\scriptsize\bfseries, text=white] at (c3) {64};
\node[dlb, below=5pt] at (c3-b) {Conv$_3$\\$3{\times}3,\;s{=}2$};
%
% Spatial MHA — 4 heads
\cub{mha}{8.0}{\cnnY}{0.9}{0.20}{0.85}{Cmha}
\node[font=\scriptsize\bfseries, text=white] at (mha) {4h};
\node[dlb, below=5pt] at (mha-b) {Spatial\\[-1pt]MHA};
%
% ---- dimension annotations along the top ----
\node[font=\tiny, text=black!40, above=1pt] at (c1-t) {$12^2$};
\node[font=\tiny, text=black!40, above=1pt] at (c2-t) {$6^2$};
\node[font=\tiny, text=black!40, above=1pt] at (c3-t) {$3^2$};
%
% ---- CNN arrows ----
\draw[arr] (gdf-r) -- (c1-l);
\draw[arr] (c1-r)  -- (c2-l);
\draw[arr] (c2-r)  -- (c3-l);
\draw[arr] (c3-r)  -- (mha-l);
%
%
% ==================== FLATTEN =============================
\node[draw, rounded corners=2pt, fill=Cconv!10, thick,
      minimum height=0.50cm, minimum width=0.65cm,
      font=\footnotesize, align=center]
    (flat) at (9.6, \cnnY) {Flat};
\node[dlb, below=2pt] at (flat.south) {576};
\draw[arr] (mha-r) -- (flat.west);
%
%
% ==================== CONCATENATION =======================
\node[draw, circle, fill=Ccat!35, thick,
      minimum size=0.55cm, inner sep=0pt, font=\normalsize]
    (cat) at (11.2, \midY) {$\oplus$};
\node[dlb, below=2pt] at (cat.south) {587};
%
%   flat (above) ↘ cat ;  sca (below) ↗ cat
\draw[arr, rounded corners=4pt] (flat.east) -| (cat);
\draw[arr, rounded corners=4pt] (sca.east)  -| (cat);
%
%
% ==================== SHARED FC TRUNK =====================
\node[draw, rounded corners=2pt, fill=Cfc!18, thick,
      minimum height=0.55cm, minimum width=1.3cm,
      font=\footnotesize, align=center]
    (sfc) at (13.0, \midY)
    {FC(256)+ReLU};
\draw[arr] (cat) -- (sfc);
%
%
% ==================== DUELING HEAD ========================
%
% ---- fork point ----
\coordinate (fork) at ($(sfc.east)+(0.4,0)$);
%
% ---- Value stream ----
\node[draw, rounded corners=3pt, fill=Cval!15, thick,
      minimum width=1.6cm, minimum height=0.65cm,
      font=\footnotesize, align=center]
    (val) at (16.0, \valY)
    {Value\\[-2pt]{\scriptsize $V(s)\!\to\!\mathbb{R}^1$}};
%
% ---- Advantage stream ----
\node[draw, rounded corners=3pt, fill=Cadv!15, thick,
      minimum width=1.6cm, minimum height=0.65cm,
      font=\footnotesize, align=center]
    (adv) at (16.0, \advY)
    {Advantage\\[-2pt]{\scriptsize $A(s,a)\!\to\!\mathbb{R}^{35}$}};
%
% ---- Y-fork arrows ----
\draw[arr, rounded corners=4pt] (sfc.east) -- (fork) |- (val.west);
\draw[arr, rounded corners=4pt] (fork) |- (adv.west);
%
%
% ==================== Q AGGREGATION =======================
\node[draw, rounded corners=3pt, fill=Cqagg!20, thick,
      minimum width=1.6cm, minimum height=0.65cm,
      font=\footnotesize, align=center]
    (qagg) at (18.8, \midY)
    {$Q{=}V{+}A{-}\overline{A}$\\[-2pt]{\scriptsize 35\;Q-values}};
%
\draw[arr, rounded corners=4pt]
    (val.east) -- ++(0.3,0) |- ($(qagg.west)+(0, 0.08)$);
\draw[arr, rounded corners=4pt]
    (adv.east) -- ++(0.3,0) |- ($(qagg.west)+(0,-0.08)$);
%
%
% ==================== ACTION MASK (7×5 GRID) ==============
\def\gX{20.8}      % grid left x
\def\gY{-0.10}     % grid bottom y  (centred ≈ \midY)
\def\gS{0.22}      % cell size
%
% ---- all cells (valid = green) ----
\foreach \r in {0,...,6}{
    \foreach \c in {0,...,4}{
        \fill[Cadv!15, draw=black!15, line width=0.15pt]
            ({\gX+\c*\gS}, {\gY+\r*\gS})
            rectangle ++({\gS-0.02}, {\gS-0.02});
    }
}
%
% ---- masked cells (red background + × mark) ----
\foreach \r/\c in {0/0,0/1,0/3,0/4, 1/0,1/4, 5/0,5/4, 6/0,6/1,6/3,6/4}{
    \fill[Cmask!12, draw=black!15, line width=0.15pt]
        ({\gX+\c*\gS}, {\gY+\r*\gS})
        rectangle ++({\gS-0.02}, {\gS-0.02});
    \draw[Cmask, line width=0.35pt]
        ({\gX+\c*\gS+0.04}, {\gY+\r*\gS+0.04})
        -- ++({\gS-0.10}, {\gS-0.10});
    \draw[Cmask, line width=0.35pt]
        ({\gX+\c*\gS+\gS-0.06}, {\gY+\r*\gS+0.04})
        -- ++({-\gS+0.10}, {\gS-0.10});
}
%
% ---- grid axis labels ----
\node[dlb, below=4pt] at ({\gX+2.5*\gS}, \gY) {$5\;a$};
\node[dlb, rotate=90]  at ({\gX-0.16}, {\gY+3.5*\gS})
    {$7\;\dot{\delta}$};
%
% ---- kinematic mask input ----
\node[draw, rounded corners=2pt, fill=Cmask!8, thick,
      minimum width=0.9cm, minimum height=0.40cm,
      font=\scriptsize, align=center]
    (kin) at ({\gX+2.5*\gS}, {\gY-0.80})
    {Kin.\ Mask\\[-2pt]{\tiny bool\;35-d}};
\draw[arr] (kin) -- ({\gX+2.5*\gS}, {\gY-0.02});
%
% ---- Q → grid ----
\draw[arr] (qagg.east) -- ({\gX-0.05}, {\gY+3.5*\gS});
%
%
% ==================== OUTPUT ==============================
\node[draw, rounded corners=3pt, fill=Cout!12, thick,
      minimum width=0.7cm, minimum height=0.55cm,
      font=\footnotesize, align=center]
    (out) at (23.5, \midY)
    {$a^*$\\[-2pt]{\scriptsize $(\dot{\delta},\,a)$}};
%
\draw[arr] ({\gX+5*\gS}, {\gY+3.5*\gS})
    -- node[above, font=\scriptsize, text=black!50] {argmax}
    (out.west);
%
%
% ==================== BACKGROUND GROUPS ===================
\begin{scope}[on background layer]
    %
    % CNN Feature Extractor
    \node[draw=Cconv!40, dashed, rounded corners=5pt,
          inner sep=7pt, fill=Cconv!3,
          fit=(c1-bl)(c1-tr)(mha-tr)(flat),
          label={[glb, Cconv!60]above:CNN Feature Extractor}] {};
    %
    % Dueling Head (MD)
    \node[draw=Cqagg!50, dashed, rounded corners=5pt,
          inner sep=7pt, fill=Cqagg!3,
          fit=(val)(adv)(qagg),
          label={[glb, Cqagg!70]above:Dueling Head (MD)}] {};
    %
    % Action Masking
    \coordinate (mask-bl) at ({\gX-0.12}, {\gY-0.12});
    \coordinate (mask-tr) at ({\gX+5*\gS+0.10}, {\gY+7*\gS+0.10});
    \node[draw=Cmask!35, dashed, rounded corners=5pt,
          inner sep=7pt, fill=Cmask!2,
          fit=(mask-bl)(mask-tr)(kin)(out),
          label={[glb, Cmask!50]above:Action Masking}] {};
    %
\end{scope}
%
\end{tikzpicture}%
}% end resizebox
%
\caption{%
Architecture of the CNN-DDQN+MD Q-network.
A convolutional branch extracts spatial features from two
$12{\times}12$ map channels
(occupancy, goal distance field)
through three convolution layers with ReLU activations
and spatial multi-head attention (MHA),
yielding a 576-dimensional feature vector.
This vector is concatenated with 11 scalar state features
(pose, velocity, steering angle, goal bearing)
and processed by a shared FC layer before splitting
into a value stream~$V(s)$ and an advantage stream~$A(s,a)$
following the dueling architecture.
The head produces 35 joint Q-values
($7$ steering rates\,$\times$\,$5$ accelerations),
and a kinematic admissibility mask zeroes infeasible actions
before argmax policy extraction.}
\label{fig:architecture}
\end{figure*}
